\section*{Hoofdstuk \ref{ch:b3:advanced-nmr} --- NMR in de diepte: pulsen, relaxatie en 2D}

\begin{solution}{exo:b3:advanced-nmr:1}
$\nu_0 = (\gamma/2\pi)B_0$: \ce{^1H} \qty{600.4}{MHz}, \ce{^{13}C} \qty{151.0}{MHz}, \ce{^{19}F}
\qty{565.1}{MHz}.
\end{solution}

\begin{solution}{exo:b3:advanced-nmr:2}
$h\nu_0/2kT = 6.626\times10^{-34} \times 400.2\times10^6/(2 \times 1.381\times10^{-23} \times 300)
= \num{3.2e-5}$. Bij \qty{4}{K} is het 75 keer groter, \num{2.4e-3} (de
benadering geldt nog).
\end{solution}

\begin{solution}{exo:b3:advanced-nmr:3}
$\theta = \gamma B_1t_p = \pi/2$ in \qty{10}{\micro s}: $\gamma B_1/2\pi = 1/(4 \times
\qty{10}{\micro s}) = \qty{25}{kHz}$. Een puls van $180^\circ$ duurt
\qty{20}{\micro s}; \qty{5}{\micro s}
geeft $45^\circ$.
\end{solution}

\begin{solution}{exo:b3:advanced-nmr:4}
$T_2 = 1/(\pi \times 3.2) = \qty{0.10}{s}$.
\end{solution}

\begin{solution}{exo:b3:advanced-nmr:5}
$(10.71/42.58)^3 = 0.0159$, maal 0.011: \num{1.75e-4}, ongeveer 1/5700 van
het protonsignaal. Omdat S/R $\propto\sqrt N$, vraagt een gelijke S/R
$5700^2 \approx \num{3e7}$ keer meer
scans --- onmogelijk; daarom gebruiken ${}^{13}$C-spectra meer
geconcentreerde monsters, ontkoppeling en polarisatieoverdracht.
\end{solution}

\begin{solution}{exo:b3:advanced-nmr:6}
$T_1 = 1.39/\ln2 = \qty{2.0}{s}$; wachttijd $\ge 5T_1 = \qty{10}{s}$.
\end{solution}

\begin{solution}{exo:b3:advanced-nmr:7}
\omterm{def:b3:advanced-nmr:experiments}{COSY}: één \omterm{def:b3:advanced-nmr:two-d}{kruispiek}, tussen het \ce{CH2}-kwartet en het \ce{CH3}-triplet van
de ethylgroep; het singlet van \ce{CH3CO} koppelt met niets. \omterm{def:b3:advanced-nmr:experiments}{HMBC} vanuit
de singletmethylgroep: het carbonylkoolstofatoom (twee bindingen) en het
\ce{CH2}-koolstofatoom (drie bindingen).
\end{solution}

\begin{solution}{exo:b3:advanced-nmr:8}
$k_c = \pi \times 50/\sqrt2 = \qty{111}{s^{-1}}$. $\Delta G^\ddagger = 8.314 \times 330 \times
\ln(6.88\times10^{12}/111) = \qty{68}{kJ/mol}$.
\end{solution}

\begin{solution}{exo:b3:advanced-nmr:9}
NOE $\propto r^{-6}$: $r_b = 0.30 \times 4^{-1/6} = \qty{0.24}{nm}$.
\end{solution}

\begin{solution}{exo:b3:advanced-nmr:10}
In het \omterm{def:b3:advanced-nmr:magnetisation}{roterende assenstelsel} blijft bij resonantie alleen $\mathbf B_1 = B_1\mathbf e_{x'}$ over:
$\dot M_{x'} = 0$, $\dot M_{y'} = \gamma B_1M_z$, $\dot M_z = -\gamma B_1M_{y'}$ (componenten van
$\gamma\mathbf M\times\mathbf B_1$). Dus $M_{x'}$ is constant en $(M_{y'}, M_z)$
draait met de hoeksnelheid $\gamma B_1$: na $t_p$ over de hoek $\gamma B_1t_p$ om $x'$.
\end{solution}

\begin{solution}{exo:b3:advanced-nmr:11}
Een spin in een iets sterker veld wint tijdens $\tau$ fase met een
constante extra snelheid; de puls van $180^\circ$ keert zijn fase om, en hij
verliest precies evenveel tijdens de tweede $\tau$: al zulke spins liggen
weer op één lijn bij $2\tau$. Toevallige spin-spinwisselwerkingen
fluctueren in de tijd en herhalen zich niet in het tweede interval: hun
defasering wordt niet ongedaan gemaakt. De echoamplitude vervalt met de
echte $T_2$.
\end{solution}

\begin{solution}{exo:b3:advanced-nmr:12}
Ethylethanoaat, \ce{CH3COOCH2CH3}: \ce{OCH2}-kwartet (4.12), \ce{CH3CO}-singlet
(2.05), \ce{CH3}-triplet; de kwartetprotonen zien het carbonylkoolstofatoom
op drie bindingen afstand, via het esterzuurstofatoom. Methylpropanoaat,
\ce{CH3CH2COOCH3}, zou zijn kwartet rond 2.3 ppm vertonen (op een
koolstofatoom rond 27 ppm in de \omterm{def:b3:advanced-nmr:experiments}{HSQC}, niet 60) en een \ce{OCH3}-singlet
rond 3.7 ppm waarvan de HMBC-piek naar het carbonyl over het zuurstofatoom
loopt.
\end{solution}

\begin{solution}{pb:b3:advanced-nmr:1}
\textbf{1.} $(2 \times 6 + 2 - 12)/2 = 1$: één C=O.
\textbf{2.} 173.6 ppm: het estercarbonyl; 60.1 ppm: de \ce{OCH2}.
\textbf{3.} Vijf; DEPT-135: \ce{CH2} negatief (60.1, 36.3, 18.6), \ce{CH3}
positief (14.3,
13.7), het carbonyl afwezig.
\textbf{4.} \qty{400.2}{MHz} en \qty{100.7}{MHz}.
\textbf{5.} $3J = \qty{21.6}{Hz} = \qty{0.054}{ppm}$.
\textbf{6.} De multipliciteiten tonen welke groepen binnen elk fragment
buren zijn, maar niet hoe de fragmenten over het zuurstofatoom en het
carbonyl met elkaar verbonden zijn.
\textbf{7.} 4.13--1.25; 2.28--1.66; 1.66--0.95 (en de symmetrische).
\textbf{8.} \ce{OCH2CH3} (4.13, 1.25) en \ce{CH2CH2CH3} (2.28, 1.66, 0.95).
\textbf{9.} De protonen van \ce{OCH2} en \ce{CH2C=O} worden gescheiden door
het zuurstofatoom en het carbonylkoolstofatoom: vijf bindingen, geen
opgeloste koppeling.
\textbf{10.} 1.66 ppm, tussen \ce{CH2} en \ce{CH3}: gekoppeld met 2 + 3 = 5
protonen met een gelijkaardige $J$, een sextet (een multiplet).
\textbf{11.} Gekoppelde methylgroepen, dus twee \ce{CH3} aan naburige
koolstofatomen (een \ce{CH3-CH3}-eenheid): hier onmogelijk.
\textbf{12.} \ce{-O-CH2-CH3} en \ce{-CH2-CH2-CH3}, plus de \ce{C=O}.
\textbf{13.} 4.13/60.1; 2.28/36.3; 1.66/18.6; 1.25/14.3; 0.95/13.7.
\textbf{14.} 14.3 (\ce{CH3} van ethyl) en 13.7 ppm (\ce{CH3} van butanoyl),
slechts \qty{0.6}{ppm} uit elkaar.
\textbf{15.} 4.13 ppm naar 173.6 ppm: H--C--O--C, drie bindingen.
\textbf{16.} 2.28 ppm (twee bindingen) en 1.66 ppm (drie bindingen) naar 173.6 ppm.
\textbf{17.} \ce{CH3CH2CH2C(=O)OCH2CH3}, ethylbutanoaat. Propylpropanoaat
zou een \ce{OCH2}-triplet (geen kwartet) rond 4.0 ppm en een kwartet rond
2.3 ppm geven.
\textbf{18.} Koppelingen over vier bindingen zijn meestal kleiner dan 1 Hz;
de HMBC-wachttijd, afgestemd op 5--10 Hz, selecteert ze niet.
\textbf{19.} Koolstofatomen relaxeren met verschillende snelheden en
krijgen door de ontkoppeling verschillende nucleaire
overhauserversterkingen.
\textbf{20.} $1 - \eu^{-5/20} = 0.22$: het carbonyl is met een factor 4
ondervertegenwoordigd.
\textbf{21.} $1 - \eu^{-5} = 0.993$.
\textbf{22.} Het NOE door protonontkoppeling, voor elk koolstofatoom
verschillend; het wordt weggewerkt door inverse-gated ontkoppeling (de
ontkoppelaar alleen aan tijdens de opname).
\textbf{23.} $256 \times \qty{100}{s} = \qty{25600}{s}$, ongeveer 7 uur.
\textbf{24.} \textbf{$5T_1$ van het carbonylkoolstofatoom: \qty{100}{s}
tussen de scans.}
\end{solution}
